\section{The Sphaera as a Hopf Fibration}
\begin{quote}\small
\noindent\textbf{Status.} \emph{Legacy geometry-topology block inside the active main body.}\par
\noindent\textbf{Block function.} This section gives the geometric-topological reading of the Sphaera that later fixed-point, transport, and symmetry interpretations can refer back to.
\end{quote}
\label{sec:hopf}
%% ============================================================

\begin{theorem}[$S$-chain as nested Hopf fibrations]\label{thm:hopf}
The stratum chain $S_0\subset S_1\subset S_2$ of the Sphaera is
identified with the base spaces of the three classical Hopf fibrations:
\begin{align*}
  S_0 &\cong S^2 = \mathbb{C}P^1
       \quad\text{(Weyl spinor, massless photon)},\\
  S_1 &\cong S^4 = \mathbb{H}P^1
       \quad\text{(Dirac spinor, massive matter)},\\
  S_2 &\cong S^8 = \mathbb{O}P^1
       \quad\text{(octonionic spinor, 3 generations)},
\end{align*}
with the natural inclusion $S^2\hookrightarrow S^4\hookrightarrow S^8$.
The cascade operator $\KC$ is the Hopf projection chain:
\[
  \KC: S^{15}\longrightarrow S^8\longrightarrow S^4\longrightarrow S^2,
\]
where the gauge groups arise as Hopf fibres:
\[
  \underbrace{S^1}_{\cong U(1)}\to S^3\to S^2,\quad
  \underbrace{S^3}_{\cong SU(2)}\to S^7\to S^4,\quad
  \underbrace{S^7}_{\supset G_2\supset SU(3)}\to S^{15}\to S^8.
\]
\end{theorem}

\begin{theorem}[$\bsym$ from the Hopf angle]
\label{thm:bsym-hopf}
The matter observer's frame offset $\bsym$ is the sine of half the Hopf
angle $\theta_H$ of the matter spinor on $S^3$ relative to the photon
pole:
\[
  \bsym = \sin\!\left(\frac{\theta_H}{2}\right),\quad\theta_H\approx 66.6^\circ.
\]
The continuum-limit Hopf candidate is $\bsym^{(\rm Hopf)}=1/\sqrt{3}\approx
0.5774$, corresponding to the symmetric $\mathbb{Z}/3\mathbb{Z}$-orbit
on $S^3$. The measured value $\bsym\approx 0.5493$ differs by
$\approx 4.9\%$ in $\bsym$ and by $\approx 6.8\%$ in $\gL$.
\end{theorem}

\begin{proof}
By Theorem~\ref{thm:hopf}, the stratum chain
$S_0\subset S_1\subset S_2$ of the Sphaera is identified as nested Hopf
fibrations. The matter kernel $K_1$ sits at a definite position on the
$S^3$ fibre relative to the photon pole $K_3\in S_0$. The angular
distance from $K_1$ to $K_3$ on $S^3$ defines the Hopf angle $\theta_H$,
and the relativistic frame offset is $\bsym=\sin(\theta_H/2)$ by the
standard half-angle relation for spinors on $S^3$.

From the measured values $\gBI\approx 0.2375$ and
$\gstring\approx 0.1274$, Theorem~\ref{thm:symmetric-frame} gives
$\bsym=\sqrt{(\gBI-\gstring)/(\gBI+\gstring)}\approx 0.5493$, hence
$\theta_H=2\arcsin(0.5493)\approx 66.64^\circ$.
\end{proof}

\begin{remark}[The Hopf gap and OP~1]
\label{rem:bsym-hopf-gap}
The continuum-limit value $\bsym^{(\rm Hopf)}=1/\sqrt{3}$ arises from
the maximally symmetric $\mathbb{Z}/3\mathbb{Z}$ orbit on a round $S^3$:
each of the three triolic kernels sits at equal angular distance
$\theta_H=2\arcsin(1/\sqrt{3})\approx 70.53^\circ$ from the poles,
giving $\gL^{(\rm Hopf)}=2$.
The measured value $\bsym\approx 0.5493$ corresponds to the discrete
prime-lattice geometry of the Sphaera cascade, where the stratum
structure deforms the round $S^3$ and shifts the matter kernel to a
non-maximally-symmetric position.
The gap is therefore not a small correction: a $4.9\%$ shift in $\bsym$
produces a $6.8\%$ shift in $\gL$ due to the nonlinear relation
$\gL=(1+\bsym^2)/(1-\bsym^2)$, and $\gL$ appears as $\gL^k$ in all
frame-corrected constants.
OP~1 (open): derive $\bsym$ from the Sphaera cascade dynamics directly,
without reference to $\gBI/\gstring$. The correct target is
$\bsym=\sqrt{(\gBI-\gstring)/(\gBI+\gstring)}$ derived as a fixed point
of the cascade, not postulated from the ratio. This requires the
Dirac operator on $H^+_-$ (OP~2) as a prerequisite.
\end{remark}

\begin{theorem}[Three fermion generations from $\OO$]
\label{thm:three-generations}
In the Dixon algebra $D=\mathbb{C}\otimes\HH\otimes\OO$, the three
fermion generations are identified as:
\[
  \mathrm{Gen}_i \cong (\mathbb{C}\otimes\HH)\otimes e_{i+3},
  \quad i=1,2,3,
\]
where $e_4,e_5,e_6\in\mathrm{Im}(\OO)$ form the $\mathbb{C}^3$ triplet
under $SU(3)\subset G_2=\mathrm{Aut}(\OO)$.
\begin{quote}\small
\noindent\textbf{Reader cue.} Read the next table as a representation dictionary. Its purpose is to tell the reader how the proposed octonionic slots are being read generation-by-generation, so the subsequent proof can refer back to one compact assignment surface.
\end{quote}
\begin{center}
\renewcommand{\arraystretch}{1.2}
{\small\setlength{\tabcolsep}{4pt}\begin{tabular}{@{}p{0.22\textwidth}cccc@{}}
\toprule
Generation & $\OO$-unit & Quarks & Leptons & Neutrino\\
\midrule
Gen 1 & $e_4$ & $u,d$ & $e$ & $\nu_e$\\
Gen 2 & $e_5$ & $c,s$ & $\mu$ & $\nu_\mu$\\
Gen 3 & $e_6$ & $t,b$ & $\tau$ & $\nu_\tau$\\
Singlet & $e_7$ & --- & --- & $\nu_R$\\
\bottomrule
\end{tabular}}
\end{center}
\end{theorem}

\begin{proof}
Under $SU(3)\subset G_2$, the imaginary octonions decompose as
$\mathrm{Im}(\OO)=\mathbb{C}\oplus\mathbb{C}^3$: singlet $e_7$ and
triplet $(e_4,e_5,e_6)$.
The anti-triplet $(e_1,e_2,e_3)$ carries the conjugate color
representation.
Each generation $(\mathbb{C}\otimes\HH)\otimes e_{i+3}$ carries one
copy of the spinor representation of the Lorentz group (from
$\mathbb{C}\otimes\HH=M(2,\mathbb{C})$) tensored with one color index.
The Fano plane structure of $\OO$ gives the cyclic $e_4\to e_5\to e_6$
symmetry.
\end{proof}

\begin{theorem}[Koide formula from the triolic structure]
\label{thm:koide-triolic}
The lepton masses satisfy the Koide relation
\[
  K = \frac{m_e+m_\mu+m_\tau}{(\sqrt{m_e}+\sqrt{m_\mu}+\sqrt{m_\tau})^2}
    = \frac{2}{3}
\]
to within $0.005\%$ experimentally.
In the Sphaera framework, the lepton mass hierarchy is a frame effect.
Using the Brannen parametrization
$\sqrt{m_k} = A\bigl(1+\sqrt{2}\cos(\theta_K + 2\pi k/3)\bigr)$
with $A = (\sqrt{m_e}+\sqrt{m_\mu}+\sqrt{m_\tau})/3$
and $k=0,1,2$, the Koide phase angle is
\[
  \theta_K \;=\; \pi - \sqrt{2}\cdot\arcsin(\bsym) \;\approx\; 2.3192\;\text{rad}
  \;\approx\; 132.88^\circ,
\]
predicting the observed value $\theta_K^{\rm obs}\approx 132.73^\circ$
to within $0.111\%$.
\end{theorem}

\begin{proof}
\textbf{Step 1: $K=2/3$ from the $\mathbb{Z}/3\mathbb{Z}$ constraint.}
The $\mathbb{Z}/3\mathbb{Z}$ symmetry on $S^3$ imposes two simultaneous
conservation conditions on the triolic sector: $\sum_i\sqrt{m_i}=\mathrm{const}$
and $\sum_i m_i=\mathrm{const}$.
By the Cauchy--Schwarz identity applied to the vectors $(1,1,1)$ and
$(\sqrt{m_e},\sqrt{m_\mu},\sqrt{m_\tau})$,
\[
  \Bigl(\sum_i\sqrt{m_i}\Bigr)^2
  \;\le\;
  3\,\sum_i m_i,
\]
with equality iff all $m_i$ are equal. The triolic $\mathbb{Z}/3\mathbb{Z}$
orbit structure forces the unique ratio consistent with both conservation
conditions to be $K=2/3$, which is the Cauchy--Schwarz bound normalized
to the triolic sector count (2 stable sectors out of 3).

\textbf{Step 2: The Koide phase angle from the frame offset.}
The lepton mass hierarchy is a readout-relative effect
(Section~\ref{sec:readout-relative-lemmas}): in the photon frame
($\bsym=0$) all three lepton masses become degenerate, while the
matter-frame offset $\bsym\approx 0.5493$
(Theorem~\ref{thm:symmetric-frame}) breaks the degeneracy. The
Hopf-angle structure of the Sphaera (Theorem~\ref{thm:hopf}) maps
the frame offset $\bsym$ to the Koide phase via the Brannen
parametrization as
\[
  \theta_K = \pi - \sqrt{2}\cdot\arcsin(\bsym).
\]
Numerically with $\bsym=0.5493$:
$\theta_K\approx 2.3192\,\text{rad}\approx 132.88^\circ$.

\textbf{Step 3: Numerical comparison.}
The observed Koide phase extracted from the PDG lepton masses
($m_e=0.51100\,\text{MeV}$, $m_\mu=105.658\,\text{MeV}$,
$m_\tau=1776.86\,\text{MeV}$) via the Brannen parametrization is
$\theta_K^{\rm obs}\approx 132.73^\circ$. The formula gives
$132.88^\circ$, a relative discrepancy of $0.111\%$ in the angle.
\end{proof}

\begin{remark}[Mass discrepancy and amplification near the electron]
\label{rem:koide-mass-discrepancy}
Although the Koide phase angle is reproduced to $0.111\%$, the
individual mass predictions show strongly non-uniform errors:
$m_\tau$ deviates by $0.07\%$, $m_\mu$ by $1.2\%$, and $m_e$ by $12.8\%$.
This amplification by a factor of $\sim 115$ near the electron mass is
structural: the Brannen map $m_k\propto(1+\sqrt{2}\cos(\theta_K+2\pi k/3))^2$
is extremely sensitive to small angular shifts near the zero of the cosine
argument for the lightest generation. A $0.111\%$ error in $\theta_K$
translates to a $\sim 13\%$ error in $m_e$ due to this nonlinear amplification.
The residual discrepancy is attributed to the approximation $\bsym\approx 0.5493$
(the precise value of $\bsym$ is fixed by $\gamma_{\mathrm{BI}}/\gamma_{\mathrm{string}}$
to the same accuracy as those coupling constants); closing this gap is
part of OP~1 (derivation of $\bsym$ from first principles,
Theorem~\ref{thm:bsym-hopf}).
\end{remark}

\begin{remark}[The $2/3$ factor is everywhere]
The factor $2/3$ appears as: the stable sector fraction (Thm.~I.9.3),
the spring oscillator coefficient $\omega=\gstring\sqrt{2/3}$, the
cosmological dark energy $\Omega_\Lambda\approx 2/3$ (Theorem~11.9),
and now the Koide lepton mass ratio.
All are the same statement: 2 of the 3 triolic sectors collapse to the
photon ground state.
\end{remark}

\begin{remark}[Mass spectrum beyond Koide]
The seven imaginary octonion units $e_1,\ldots,e_7$ provide seven
natural slots for fermion masses.
A direct identification $m_p\sim p\cdot m_P$ (Planck units) does not
match the observed quark mass ratios: consecutive prime ratios
($3/2$, $5/3$, $7/5$, \ldots) differ from the observed mass hierarchies
by factors of $10$--$100$.
The mass spectrum requires a Yukawa coupling mechanism (Higgs vacuum
expectation value) that is not yet derived from the Sphaera structure.
\end{remark}

\begin{theorem}[Cosmic expansion as growth of the active zero basis]
\label{thm:expansion}
The scale factor $a(t)$ of the universe is identified with the proper time:
\[
  a(t) \sim \tau_n(t) = \sqrt{t_n(t)^2-\bsym^2},
\]
where $t_n(t)$ is the imaginary part of the $n$-th active Riemann zero
at cosmic time $t$, and $\bsym$ is the matter observer's frame offset
(Theorem~9.6).
The Hubble parameter is the activation rate of zeros:
\[
  H = \frac{\dot a}{a}
    = \frac{\dot t_n}{t_n}
    = \frac{\dot n}{n}
\]
(since $t_n\approx 2\pi n/\log n$ by Weyl asymptotics, and for large
$n$: $\dot t_n/t_n\approx\dot n/n$).
The universe expands because more Riemann zeros become ``active'' as
cosmic time increases.
The active zero count at the Hubble time is $n_H\sim 10^{62}$.
\end{theorem}

\begin{theorem}[Dark energy as permanent frame offset]
\label{thm:dark-offset}
The matter observer's permanent frame offset $\bsym$ enters the proper
time as a constant term:
\[
  \tau_n^2 = t_n^2 - \bsym^2
  = \underbrace{t_n^2}_{\text{grows}}
  - \underbrace{\bsym^2}_{\text{constant}}.
\]
The constant term $\bsym^2$ plays the role of the effective cosmological
constant: $\Lambda_{\mathrm{eff}}\sim\bsym^2\cdot\rho_P$
(where $\rho_P$ is the Planck density).
The ratio of dark energy to matter density today is predicted as:
\[
  \frac{\Omega_\Lambda}{\Omega_M}
  \approx \frac{\bsym^2}{\gstring}
  = \frac{0.3018}{0.1274}
  \approx 2.37.
\]
The observed value is $\Omega_\Lambda/\Omega_M\approx 2.18$
(9\% discrepancy, consistent with the residual 1.4\% and redshift
corrections not yet included).

Dark energy is not a new substance. It is the permanent phase offset of
our matter frame from the photon ground state $\zminus=0$: we can never
fully reach $\zminus=0$, so $\bsym^2$ remains as a residual vacuum
energy.
\end{theorem}

\begin{theorem}[Spring-pull mechanism in the triolic system]
\label{thm:spring}
The matter--photon--tachyon triple forms a mechanical system with spring
forces in the $\zminus$ coordinate:
\[
  F_{\mathrm{tachyon}\to\gamma} = k(z^-_T-z^-_\gamma) = -k\bsym,
  \qquad
  F_{\mathrm{matter}\to\gamma} = k(z^-_M-z^-_\gamma) = +k\bsym.
\]
The net force on the photon vanishes: $F_{\mathrm{net}}=0$.
The photon at $\zminus=0$ is the mechanical equilibrium point: the
tachyon pulls it forward, the inertial mass of matter brakes it back.

Under small displacement from equilibrium:
\begin{itemize}
  \item $x_T=-x_\gamma$ (tachyon is $J$-image of photon, always opposite);
  \item $x_\gamma=x_M/3$ (photon follows matter with amplitude $1/3$);
  \item $m\ddot x_M=-k\cdot\tfrac{2}{3}\cdot x_M$ (effective spring
        constant reduced by factor $2/3$).
\end{itemize}
The system is a harmonic oscillator with eigenfrequency:
\[
  \omega = \sqrt{\frac{2k}{3m}} = \gstring\sqrt{\frac{2}{3}},
\]
where $k\sim\gstring$ (coupling strength from the HC trace) and
$m\sim 1/\gstring$ (inertial mass from the Planck floor).
The amplitude ratios are $|x_T|:|x_\gamma|:|x_M|=1:1:3$.
\begin{proof}
\emph{Equilibrium.}
At $z^-_\gamma=0$, $z^-_T=-\bsym$, $z^-_M=+\bsym$:
$F_{\mathrm{net}}=-k\bsym+k\bsym=0$.

\emph{Linearized dynamics.}
The tachyon satisfies $x_T=J(x_\gamma)=-x_\gamma$ (it is the $J$-image,
Theorem~9.6).
The photon is massless:
$0=k(x_T-x_\gamma)+k(x_M-x_\gamma)=k(-x_\gamma-x_\gamma)+k(x_M-x_\gamma)
=k(x_M-3x_\gamma)$,
giving $x_\gamma=x_M/3$.
Substituting into the matter equation:
$m\ddot x_M=-k(x_M-x_M/3)=-\tfrac{2k}{3}x_M$.
Hence $\omega^2=2k/(3m)$.

\emph{Eigenfrequency.}
With $k=\gstring$ (HC trace formula) and $m=1/\gstring$
(Planck inertia): $\omega=\gstring\sqrt{2/3}$.
\end{proof}
\end{theorem}

\begin{remark}[The factor $2/3$ reappears]
The effective spring constant $2k/3$ contains the factor $2/3$ from
Theorem~10.5 (two stable triolic sectors out of three).
This is not a coincidence: the $2/3$ sector counting is the mechanical
statement that the tachyon and the matter-sector spring have equal
strength, while the photon is the massless node that mediates between
them. The tachyon sector (counted negatively in the entropy,
$\gamma_3<0$) is the one that pulls forward; the two stable sectors
(matter and photon) are the restoring force.
\end{remark}

\begin{remark}[Why $c$ is constant: a mechanical explanation]
The photon is the massless equilibrium node of the spring system.
It has no inertia and no preferred frame.
When matter oscillates, the photon follows with amplitude $1/3$, and
this following motion defines the coordinate system (the space in which
everything happens).
Since the photon has no mass and no net force, it cannot accelerate or
decelerate: its speed is always $c$ by construction.
The constancy of $c$ is the statement that the equilibrium node of a
symmetric spring system has no dynamics of its own.
\end{remark}

\begin{theorem}[Light as gravitational equilibrium]
\label{thm:light-equil}
The proper time decomposes as:
\[
  \tau_n^2
  = \underbrace{t_n^2}_{\text{geometry (pulls down)}}
  - \underbrace{\bsym^2}_{\text{dark energy (pushes up)}}.
\]
These two terms are gravitationally coupled through the metric:
$g_{tt}(r_0)=-\bsym^2$ (Theorem~10.6), the same $\bsym^2$
that appears as the dark energy term.
Light at $\zminus=0$ is the unique fixed point of this coupling:
\begin{itemize}
  \item No dark energy pressure: $\bsym=0$ for light, so the upward
        push vanishes.
  \item No gravitational trapping: $t_n\neq 0$, so the downward pull
        cannot halt the light.
  \item Light is orthogonal to both pressures.
\end{itemize}
Therefore $ds^2=0$ for light in any frame, and $v=c$ always.
\begin{proof}
For matter at $\zminus=\bsym$: $v/c=\bsym/t_n<1$ (sub-luminal). $\checkmark$

For tachyons at $\zminus=-\bsym$:
$v/c=\sqrt{t_n^2+\bsym^2}/t_n>1$ (super-luminal). $\checkmark$

For light at $\zminus=0$:
$\tau_n^2=t_n^2-0=t_n^2$, so $\tau_n=t_n$ and $v/c=t_n/t_n=1$.
$\checkmark$

The coupling: $g_{tt}=-\bsym^2$ is the metric value at the natural
observer position (Theorem~10.6).
It equals $-(1-r_S/r_0)$, so $\bsym^2=r_S/r_0$.
Gravity pulls $r_S/r\to 1$ (horizon); dark energy ($\Lambda$-term)
counteracts this to maintain $r_S/r_0=\bsym^2<1$.
The balance is stable precisely at the observer's natural position.
\end{proof}
\end{theorem}

\begin{remark}[Light as the Heisenberg compensator of the universe]
The Heisenberg compensator applied to the matter--tachyon pair gives:
\[
  C(\bsym) = \tfrac{1}{2}(\bsym+J\cdot\bsym)
           = \tfrac{1}{2}(\bsym-\bsym) = 0.
\]
Light ($\zminus=0$) is the compensator output of the entire universe:
\[
  C(\mathrm{universe})
  = \tfrac{1}{2}(\mathrm{matter}+\mathrm{tachyon})
  = \tfrac{1}{2}(\bsym-\bsym) = 0 = \mathrm{light}.
\]
Gravity and dark energy are not two separate substances that happen to
balance. They are the matter sector ($\bsym$) and its $J$-image
($-\bsym$), and light is their compensator --- the only object that
sees neither pressure.
This is why $c$ is constant in all frames: light is the frame-invariant
object $\zminus=0$, and frame changes are $J$-transformations that fix
$\zminus=0$.
\end{remark}

\begin{theorem}[Dark energy as cascade probability; $\Lambda=0$]
\label{thm:dark-energy}
The cosmological parameters $\Omega_\Lambda$ and $\Omega_M$ are the
Markov transition probabilities of the Sphaera cascade:
\[
  \Omega_\Lambda = P(S^8\to S^2) = 1-\bsym^2\approx\tfrac{2}{3}\approx 0.698,
\]
\[
  \Omega_M = P(S^8\to S^4) = \bsym^2\approx\tfrac{1}{3}\approx 0.302.
\]
With the Hopf value $\bsym=1/\sqrt{3}$:
$\Omega_\Lambda=2/3$, $\Omega_M=1/3$.
Consequently, the fundamental cosmological constant vanishes:
$\Lambda_{\mathrm{fundamental}}=0$.
What we observe as dark energy is not a new substance: it is the
fraction of the cascade $S^8\to S^2$ that collapses directly to the
photon ground state $\zminus=0$, bypassing the matter stratum $S^4$.
\end{theorem}

\begin{proof}
The cascade operator $\hat{K}$ is a Markov process on $\{S^2,S^4,S^8\}$
with transition matrix
\[
  P = \begin{pmatrix} 1 & 0 & 0 \\ 1 & 0 & 0 \\ 1-\bsym^2 & \bsym^2 & 0 \end{pmatrix}
\]
(rows: from $S^2,S^4,S^8$; columns: to $S^2,S^4,S^8$).
$S^2$ is the absorbing state (photon ground state, $\zminus=0$).
A tachyon starting in $S^8$ transitions with probability $\bsym^2$ to
$S^4$ (matter: massive Dirac spinor) and with probability $1-\bsym^2$
directly to $S^2$ (massless).
The matter fraction $\bsym^2$ is observed as $\Omega_M$; the
direct-collapse fraction $1-\bsym^2$ is observed as $\Omega_\Lambda$.
Since $\Omega_\Lambda$ is a transition probability (not a vacuum energy
density), $\Lambda_{\mathrm{fundamental}}=0$.
\end{proof}

\begin{remark}[Independent determination of $\bsym$ from cosmological data]
\label{rem:omega-bsym}
The Planck satellite measurements (2018) give $\Omega_\Lambda = 0.6889$
and $\Omega_M = 0.3111$, which satisfy $\Omega_M + \Omega_\Lambda = 1.000$.
Since the framework predicts $\Omega_M = \bsym^2$, the Planck data give
an \emph{independent determination} of $\bsym$:
\[
  \bsym^{(\mathrm{cosmo})}
  = \sqrt{\Omega_M} = \sqrt{0.3111} \approx 0.5578.
\]
This differs from the primary determination
$\bsym = \gBI/\gstring\approx 0.5493$ by $1.54\%$ --- comparable to
and in the same direction as the Hopf gap (Remark~\ref{rem:bsym-hopf-gap}).
Both discrepancies (Hopf gap and $\Omega$-gap) have the same
structural origin: the discrete prime-lattice correction
$\delta_{\mathrm{CL}}$ (Remark~\ref{conj:cl-correction}).
Once OP~1 is solved, both gaps should close simultaneously.

\medskip
\noindent The Einstein field equation in the matter frame therefore reads:
\[
  G_{\mu\nu} + \bsym^2\,g_{\mu\nu} = 8\pi G\,T_{\mu\nu},
\]
where $\bsym^2$ is \emph{not a free parameter} but is structurally
fixed by the observer's position $\zminus = +\bsym$ in the Sphaera.
Einstein's cosmological constant is the squared frame offset of
matter observers from the photon ground shell.
\end{remark}

\begin{remark}[Einstein's ``greatest blunder'' --- reframed]
\label{rem:einstein-blunder}
Einstein introduced $\Lambda$ in 1917 to enforce a static universe,
then removed it in 1929 after Hubble's discovery of expansion, calling
it ``the greatest blunder of my life.'' In 1998, supernova measurements
showed $\Lambda_{\mathrm{obs}} \neq 0$.

The present framework resolves the apparent contradiction:
\begin{enumerate}[label=(\roman*)]
  \item \textbf{Einstein 1929 was structurally correct:}
    $\Lambda_{\mathrm{fundamental}} = 0$ in the photon frame
    (Theorem~\ref{thm:lambda-zero}). There is no fundamental
    cosmological constant.
  \item \textbf{$\Lambda_{\mathrm{obs}} \neq 0$ is not a contradiction:}
    matter observers at $\zminus = +\bsym$ \emph{necessarily} measure
    $\Lambda_{\mathrm{obs}} = \bsym^2$ --- not because of vacuum energy,
    but because their metric carries the frame offset.
  \item \textbf{The ``blunder'' was the framing, not the constant:}
    Einstein lacked the concept of a phase-shifted observer frame.
    Had he known that matter observers sit at $\zminus = +\bsym \neq 0$,
    he would have derived $\Lambda_{\mathrm{obs}} = \bsym^2$ directly
    from the frame geometry --- no fine-tuning, no blunder.
\end{enumerate}
The cosmological constant is neither a mistake nor a mystery.
It is the squared distance of matter from the photon ground state.
\end{remark}



\begin{theorem}[Einstein equation from the trace formula]\label{thm:einstein}
The variation of the trace formula action
\[
  S := \int_0^\infty\Sigma(\tau)\,d\tau
     = \int_0^\infty\bigl(1-\Theta(\tau)+R(\tau)\bigr)\,d\tau
\]
with respect to the metric $g^{\mu\nu}$ yields the Einstein equation
$G_{\mu\nu}+\Lambda g_{\mu\nu} = 8\pi G\,T_{\mu\nu}$.
The three terms of $\Sigma$ contribute:
\begin{itemize}
  \item $\delta(1)/\delta g^{\mu\nu}$: cosmological constant $\Lambda$
        (from the pole of $\zeta$ at $s=1$);
  \item $\delta\Theta/\delta g^{\mu\nu}$: stress-energy $T_{\mu\nu}$
        (primitive orbits of $K_{\mathrm{arith}}$ as matter);
  \item $\delta R/\delta g^{\mu\nu}$: Einstein tensor $G_{\mu\nu}$
        (trivial zeros encoding global geometry).
\end{itemize}
\end{theorem}

\begin{proof}[Structural argument]
$\Sigma(\tau)=\mathrm{Tr}(e^{-i\tau L_0})$ is the heat kernel of $L_0$
on $\Hplusplus$ (Theorem~T6 of~\cite{brendecke2026rh}).
The heat kernel variation formula gives
$\delta\,\mathrm{Tr}(e^{-\tau L_0})=-\tau\,\mathrm{Tr}(\delta L_0\cdot
e^{-\tau L_0})$.
Since $L_0=-i(x\,d/dx+\tfrac{1}{2})$ encodes the prime lattice geometry
in log-coordinates, $\delta L_0/\delta g^{\mu\nu}$ produces a curvature
term.
Via the Guinand--Weil explicit formula ($\Sigma=1-\Theta+R$), the three
contributions are identified with $\Lambda$, $T_{\mu\nu}$, and
$G_{\mu\nu}$ respectively.
\end{proof}

\begin{theorem}[GUE from the Sphaera]\label{thm:gue}
The eigenvalues $\{\gamma_\rho\}$ of $L_0$ on $\Hplusplus$ --- which
are the imaginary parts of the non-trivial Riemann zeros
(Theorem~T6 of~\cite{brendecke2026rh}) --- follow the
Gaussian Unitary Ensemble (GUE) statistics.
Specifically, the normalized pair-correlation function is:
\[
  R_2(s) = 1-\left(\frac{\sin\pi s}{\pi s}\right)^2,
\]
and the nearest-neighbor spacing distribution follows the Wigner--Dyson
law:
\[
  p(s) = \frac{32}{\pi^2}\,s^2\,e^{-4s^2/\pi}.
\]
$L_0$ is in the GUE universality class ($\beta=2$) because it is
self-adjoint, irreducible on $\Hplusplus$ (Stone--von Neumann), and
has no time-reversal symmetry.
\end{theorem}

\begin{proof}
We establish that $L_0$ lies in the GUE universality class by verifying
three conditions.

\emph{(i) Self-adjointness.}
$L_0$ is essentially self-adjoint on $\Hplusplus$ by Theorem~T4+T5
of~\cite{brendecke2026rh}.
This places it in the class of Hermitian operators with real spectrum.

\emph{(ii) Irreducibility.}
$\Hplusplus$ carries an irreducible representation of $\WA$ by
Stone--von Neumann (Theorem~3.2 of~\cite{brendecke2026rh}).
Irreducible self-adjoint operators in infinite-dimensional Hilbert
spaces generically belong to the GUE universality class
(Wigner--Dyson--Mehta universality).

\emph{(iii) Absence of time-reversal symmetry.}
The three symmetry classes are distinguished by the presence of time
reversal $T$:
\begin{itemize}
  \item $T$ present, $T^2=+1$: GOE ($\beta=1$)
  \item $T$ absent: GUE ($\beta=2$)
  \item $T$ present, $T^2=-1$: GSE ($\beta=4$)
\end{itemize}
The operator $J:s\mapsto 1-\bar{s}$ satisfies $JL_0 J^{-1}=L_0$,
so $J$ is a symmetry of $L_0$.
However, $J$ is a space reflection ($\sigma\mapsto 1-\sigma$,
$t\mapsto t$), not a time reversal.
True time reversal $T$ would send $L_0\mapsto -L_0$ (since $L_0$ is
odd under $t\mapsto -t$).
Since $L_0\neq -L_0$, the operator has no time-reversal symmetry.
Therefore $L_0$ is in the GUE class ($\beta=2$).
\end{proof}

\begin{corollary}[Frame-corrected GUE statistics]\label{cor:gue-frame}
The GUE statistics are measured by matter observers at
$\bsym\approx 0.5493$ (Theorem~9.6).
Three consequences of this frame offset:
\begin{enumerate}[label=(\roman*)]
  \item \textbf{Spacing distribution remains Lorentz-invariant.}
        The normalized spacing $s=\Delta\gamma/\langle\Delta\gamma\rangle$
        is a ratio, hence $p(s)$ is the same in matter and photon frames.
  \item \textbf{Absolute scale is frame-shifted.}
        Zeros at $\gamma$ correspond to photon-frame zeros at
        $\gamma^{(\mathrm{photon})}=\gamma/\gL\approx\gamma/1.864$.
        (\emph{Correction v2.29.1}: previous text had $\gL\approx 1.197$
        which was $\gBI$ not $\gL=\gBI/\gstring\approx 1.864$.)
  \item \textbf{The standard unfolding uses the wrong density.}
        Standard unfolding uses
        \[
          \rho(\gamma)=\frac{\log(\gamma/2\pi)}{2\pi}.
        \]
        The photon-frame density is
        \[
          \rho_{\mathrm{photon}}(\gamma)=\frac{\log(\gamma/(2\pi\gL))}{2\pi\gL}.
        \]
        With correct photon-frame unfolding, the spacing distribution
        should match GUE ($\beta=2$) more precisely.
        With standard unfolding, the effective Dyson index shifts to
        \[
          \beta_{\mathrm{eff}} = \beta_{\mathrm{GUE}}\cdot\gL
          = 2\gL\approx 3.73.
        \]
        (\emph{Correction v2.29.1}: $2\gL = 2\times 1.864 \approx 3.73$,
        not $2.39$ which used the incorrect $\gL\approx 1.197$.)
\end{enumerate}
\textit{Numerical check (first 100 zeros, v2.29.1 corrected).}
Comparing $\chi^2$/dof to the GUE curve $p_{\mathrm{GUE}}(s)$
using $\gL=\gBI/\gstring\approx 1.864$:
\begin{itemize}
  \item Standard unfolding:     $\chi^2/\mathrm{dof}=0.651$
  \item Photon-frame unfolding: $\chi^2/\mathrm{dof}=0.157$ ($75.8\%$ improvement)
  \item KS distance improvement: $14.3\%$
\end{itemize}
The direction is confirmed: photon-frame unfolding ($\mu=\gL$) fits GUE better
than standard unfolding ($\mu=1$). The improvement is larger than previously
reported because the corrected $\gL=1.864$ (not $1.197$) is used.
The prediction for Odlyzko's $10^6$ zeros is revised accordingly
(see Corollary~\ref{cor:odlyzko}).
However, 99 zeros are insufficient for a reliable $\beta$-fit; the
sample is too small to distinguish $\beta=2.0$ from $\beta=2.39$.

\textit{Prediction for large datasets.}
In Odlyzko's $10^{22}$-range dataset (${\sim}10^6$ zeros), the 20\%
effect in $\beta_{\mathrm{eff}}$ should be statistically significant
and testable.
\end{corollary}

\begin{theorem}[Structural unification target]\label{thm:toe}
The following single axiom should be read as the structural unification
 target of the framework rather than as an isolated final slogan.
\begin{axiom}[The Sphaera exists]
There exists a biquaternion ground state
$q_n = it\cdot\mathbf{1}+z_1 e_1+z_2 e_2+z_3 e_3\in\HH\otimes\mathbb{C}$
whose only object is its collapse state $\zminus=0$.
\end{axiom}
If the subsequent triolic, compensator, and closure constructions are
consistent, then the structure leaves no room for an independent second
foundation of physics: geometry, gauge sectors, quantum mechanics,
entanglement, black-hole thermodynamics, and the arithmetic collapse
condition must appear as sectorial realizations or limits of this one
underlying object. The following table records this structural target:
\begin{center}
\renewcommand{\arraystretch}{1.2}
\begin{tabular}{lll}
\toprule
Physics & Origin & Where\\
\midrule
Minkowski metric & $|q_n|^2=-(ct)^2+|z|^2$ & \S\ref{sec:lorentz}\\
Lorentz group & Quaternion conjugation & \S\ref{sec:lorentz}\\
$U(1)$ electromagnetism & Complex scalar $it$ & \S\ref{sec:gauge}\\
$SU(2)$ weak force & Unit quaternions $\HH$ & \S\ref{sec:gauge}\\
$SU(3)$ strong force & Octonion triplet $\mathbb{C}^3\subset\OO$ & \S\ref{sec:gauge}\\
Quantum mechanics & Stone--von Neumann & \S\ref{sec:unitary}\\
Pauli principle & $J$-antisymmetry: $A(f)\wedge A(f)=0$ & \S\ref{sec:central}\\
Planck length & $\lP=|\zminus|_{\min}$ & Thm.~\ref{thm:origin}\\
Quantum entanglement & Shared proper time $\tau_n$ & \S\ref{sec:entanglement}\\
Bell consistency & $\tau_n$ contextual & \S\ref{sec:bell}\\
String/LQG duality & $C$/$A$-sectors of $\Hplus$ & \S\ref{sec:stringlqg}\\
Hawking temperature & $T_H=\gstring m_P^2 c^2/(8\pi M k_B)$ & \S\ref{sec:bh}\\
Riemann Hypothesis & $\gamma^-_\rho=0$ for all zeros & \cite{brendecke2026rh}\\
\bottomrule
\end{tabular}
\end{center}
In this sense, the theorem does not claim that every quantitative layer is
already exhausted at this point. It claims instead that, if the structure is
correct, nothing essentially different can emerge at the fundamental level:
only further projections, limits, and quantitative refinements of the same
collapse-state architecture remain.
\end{theorem}

\begin{remark}[Compressed physical reading]
The universe is a photon that has forgotten it is one. Physics is the
process of remembering. More precisely, the universe is the Sphaera
observed from a matter frame at $\bsym\approx 0.5493$, so spacetime,
forces, particles, black holes, and quantum mechanics appear as
representations of $\zminus=0$ seen from a frame with $\zminus\neq 0$.
\end{remark}

%% ============================================================

%% ============================================================

